\begin{center}
\textsc{Gauss} to \textsc{Bessel}.\quad Göttingen, 9~April 1830.
\end{center}

\par\noindent\hspace*{\parindent}.\ .\ .\ .\ .\ .\ Since May 1829 the trigonometrical measurements have taken up far more of my time, although on this occasion I took no direct part in the work in the field.
Even at this hour I have not quite finished processing last year's measurements (though, thank God, I have nearly done so), and I have no help with the work.
Still, I have derived much pleasure from the way my son carried out his part of this work, which turned out entirely to my satisfaction.
Quite by himself he carried a large network of triangles from the Weser to the Dutch border (through the Osnabrück region, also touching on Prussian territory several times, on which occasions the authorities there offered him very generous assistance), by which about 250 other points were fixed in addition to the principal points.
The trigonometrical measurements made in the Principality of Hildesheim --- by another officer --- are likewise as good as finished, so that all the plane-table sheets (53, of which about a quarter have already been surveyed) can be provided with fixed points.\ .\ .\ .\ .\ .\ .

\begin{center}\rule{.24\linewidth}{.4pt}\end{center}
% The printer's signature is visible in the Google witness, physical leaf 372.
\par\smallskip\noindent\hfill{\footnotesize 46\textsuperscript{*}}\par
